\documentclass[10pt]{article}

\usepackage[margin=0.5in]{geometry}
\usepackage{amsmath,amssymb}
\usepackage[T1]{fontenc}
\usepackage{lmodern}
\usepackage{xcolor}
\usepackage[hidelinks]{hyperref}

\title{What \texttt{fit\_exp\_cnn.py} does}
\author{}
\date{}

\begin{document}

\maketitle

The script gives a small numerical test of the Cartan, or group-valued,
network layers in this repository.  The target is the scalar function
\(f(x)=e^x\) on the interval \([0,1]\).  We first construct the training set
\begin{equation}
 x_i=\frac{i}{255},\quad i=0,\ldots,255,
 \qquad y_i=e^{x_i},
\end{equation}
and then minimize the mean-square error between the network output and \(y_i\).
The training step is full-batch: all 256 samples are used in every epoch.

\section{Representation and network}

The input is first embedded into the Cartan--fiber representation used by the
code.  For a scalar input this is simply
\begin{equation}
 x \longmapsto (c,f)=(0,x),
\end{equation}
where \(c\) denotes the Cartan coordinate and \(f\) the fiber coordinate.
The pair \((c,f)=(0,x)\) in Eq.~(2) is therefore the input to
\texttt{HyperbolicLinear}(2,32) in Eq.~(3).  The first component is the Cartan
coordinate and the second is the fiber coordinate; the layer maps this
two-component representation to a 32-component Cartan--fiber representation.

\noindent\fcolorbox{blue}{blue!5}{\parbox{0.94\linewidth}{\scriptsize
\textbf{\textcolor{blue}{Definition.}} \texttt{HyperbolicLinear}(n,m), denoted
\(\mathrm{HLinear}_{n\to m}\), maps an \(n\)-component representation (one
Cartan and \(n-1\) fiber coordinates) to an \(m\)-component one.  Its ordinary
fiber weight matrix has size
\((m-1)\times(n-1)\).  \texttt{HyperbolicRegression}(n,1), denoted
\(\mathrm{HReg}_{n\to1}\), maps \(\mathfrak{s}_n\) to one ordinary scalar.}}

The model then applies
\begin{equation}
\begin{split}
 &\texttt{HyperbolicLinear}(2,32),\quad
   \texttt{HyperbolicActivation}(\mathrm{Id}),\\
 &\texttt{HyperbolicLinear}(32,32),\quad
   \texttt{HyperbolicActivation}(\mathrm{Id}),
\end{split}
\end{equation}
\begin{equation}
\scriptstyle
\begin{aligned}
 \mathfrak{s}_{2}
 &\xrightarrow{\mathrm{HLinear}_{2\to32}} \mathfrak{s}_{32}
  \xrightarrow{\mathrm{HAct}_{\mathrm{Id}}} \mathfrak{s}_{32}
  \xrightarrow{\mathrm{HLinear}_{32\to32}} \mathfrak{s}_{32} \\
 &\xrightarrow{\mathrm{HAct}_{\mathrm{Id}}} \mathfrak{s}_{32}
  \xrightarrow{\mathrm{HReg}_{32\to1}} \mathbb{R}.
\end{aligned}
\end{equation}
The two activation layers
preserve the Cartan coordinate and apply the requested pointwise activation to
the fiber.  Since the script passes \(\mathrm{Id}\), they introduce no
additional pointwise nonlinearity; the algebra-aware linear maps and their
group operations still determine the representation used by the regression
layer.

\noindent\fcolorbox{orange}{yellow!15}{\parbox{0.94\linewidth}{\scriptsize
\textbf{Where is the exponential map?}  No explicit manifold exponential map
is used in the forward pass: \(e^{c}\) in the regression is an ordinary
scalar exponential.  Manifold retractions may be used by the optimizer.}}

Schematically, a hyperbolic linear layer first forms a fiber linear map,
\(f\mapsto Wf+b\), while retaining the Cartan coordinate.  It can then apply
a group-valued bias and a fiber rotation.  The final regression maps the
resulting representation to an ordinary scalar.  In the implementation its
main terms have the form
\begin{equation}
 \mathcal R(c,f)=e^c W_Rf+(1+\|f\|^2)e^c\alpha+e^{-c}b_R,
\end{equation}
which is the scalar expression evaluated by \texttt{hregression} in the
repository.  The notation distinguishes the regression weight \(W_R\) and
bias \(b_R\) from the parameters of the preceding linear layers.

The complete forward map, including the matrix and group operations, is
conveniently written as

\begin{equation}
\begin{aligned}
 z_0&=(c_0,f_0)=(0,x),
 &z_1&=\mathcal L_1(z_0),\\
 z_2&=\mathcal A(z_1),
 &z_3&=\mathcal L_2(z_2),\\
 z_4&=\mathcal A(z_3),
 &\widehat f_\theta(x)&=\mathcal R(z_4),
 \end{aligned}
\qquad
\begin{aligned}
 \mathcal L_i(c,f)&=R_{\theta_i}\!\left(\beta_i*(c,W_i f+b_i)\right),\\
 \beta*(c,g)&=(c+\beta_c,\,g+e^{-c}\beta_f),\\
 \mathcal A(c,f)&=(c,f),\\
 \mathcal R(c,f)&=e^c W_Rf+(1+\|f\|^2)e^c\alpha+e^{-c}b_R.
 \end{aligned}
\end{equation}

Here \(W_1\in\mathbb R^{31\times1}\), \(W_2\in\mathbb
R^{31\times31}\), and \(W_R\in\mathbb R^{1\times31}\); the maps
\(R_{\theta_i}\) are the learned fiber rotations.

The trainable network parameters are
\begin{equation}
 \Theta=\{W_1,b_1,\beta_1,\theta_1,
          W_2,b_2,\beta_2,\theta_2,
          W_R,\alpha,b_R\}.
\end{equation}
Here \(b_i\in\mathbb R^{31}\), while \(\beta_i\) and \(\theta_i\) each have
32 components; \(\alpha,b_R\in\mathbb R\).  The embedding and the identity
activations have no trainable parameters.  Counting these entries gives 1215
trainable scalar parameters, before distinguishing their Euclidean and
manifold constraints.

The learned fiber rotation \(R_{\theta_i}\) is an isometry acting on the
Cartan--fiber representation after the matrix map and group-valued bias.  Its
parameter \(\theta_i\) mixes the Cartan and fiber coordinates without changing
the representation dimension.  It is optimized, together with the other
parameters, by \texttt{RiemannianAdam}; it is therefore a geometric layer
operation, not an additional pointwise activation.

\section{Optimization and check}

The random seed is fixed with \texttt{torch.manual\_seed(0)}.  With the
default settings, the network is trained for 3000 epochs with learning rate
\(10^{-2}\).  The loss is
\begin{equation}
 \mathcal L=\frac{1}{256}\sum_{i=1}^{256}
 \bigl(\widehat f(x_i)-e^{x_i}\bigr)^2.
\end{equation}
This is a deterministic approximation test rather than a generalization
experiment: the reported error is evaluated on a separate, denser plotting
grid, but no held-out dataset is used during optimization.

The script uses \texttt{RiemannianAdam}, because some layer parameters live on
manifolds rather than in an unconstrained Euclidean parameter space.  Each
epoch computes the prediction, backpropagates the loss, and performs one
manifold-aware optimizer step.

After training, the model is evaluated on 500 points.  The script reports
\begin{equation}
 \max_{x\in\text{plot grid}}|\widehat f(x)-e^x|
\end{equation}
and plots the exact curve together with the learned approximation; with
\texttt{--save exp\_fit.png}, it is written to a file.

\end{document}
